%====================================================%
% SEZIONE 1: Dominio di una funzione
%====================================================%

\begin{question}[ID=1.1, subtitle={Dominio di una funzione}, difficulty=facile, tags={funzioni, dominio}]
Qual è il dominio della funzione $f(x)=\frac{1}{x-3}$?

\begin{tabularx}{\linewidth}{@{}X X X X@{}}  
\boxed{A}\hspace{6pt} $\mathbb{R}$ &  
\boxed{B}\hspace{6pt} $\mathbb{R}\setminus\{3\}$ &  
\boxed{C}\hspace{6pt} $[3,+\infty)$ &  
\boxed{D}\hspace{6pt} $\mathbb{R}\setminus\{0\}$  
\end{tabularx}

\begin{solution}
\begin{enumerate}
    \item La risposta corretta è $B)$, perché il denominatore deve essere diverso da zero: $x\neq3$.
\end{enumerate}
\end{solution}
\end{question}

%----------------------------------------------------%

\begin{question}[ID=1.2, subtitle={Dominio di una funzione irrazionale}, difficulty=facile, tags={funzioni, dominio}]
Qual è il dominio della funzione $f(x)=\sqrt{x+2}$?

\begin{tabularx}{\linewidth}{@{}X X X X@{}}  
\boxed{A}\hspace{6pt} $x>-2$ &  
\boxed{B}\hspace{6pt} $x\geq-2$ &  
\boxed{C}\hspace{6pt} $x\leq-2$ &  
\boxed{D}\hspace{6pt} $\mathbb{R}$  
\end{tabularx}

\begin{solution}
\begin{enumerate}
    \item La risposta corretta è $B)$, perché il radicando deve essere maggiore o uguale a zero: $x+2\geq0$.
\end{enumerate}
\end{solution}
\end{question}

%----------------------------------------------------%

\begin{question}[ID=1.3, subtitle={Dominio di una funzione razionale e irrazionale}, difficulty=media, tags={funzioni, dominio}]
Determina il dominio della funzione $f(x)=\frac{\sqrt{x-1}}{x-2}$.

\begin{tabularx}{\linewidth}{@{}X X X X@{}}  
\boxed{A}\hspace{6pt} $[1,+\infty)$ &  
\boxed{B}\hspace{6pt} $[1,+\infty)\setminus\{2\}$ &  
\boxed{C}\hspace{6pt} $(-\infty,1]\setminus\{2\}$ &  
\boxed{D}\hspace{6pt} $\mathbb{R}\setminus\{2\}$  
\end{tabularx}

\begin{solution}
\begin{enumerate}
    \item La risposta corretta è $B)$, perché $x\geq1$ per la radice e $x\neq2$ per il denominatore.
\end{enumerate}
\end{solution}
\end{question}

%----------------------------------------------------%

\begin{question}[ID=1.4, subtitle={Dominio di una funzione con radice e frazione}, difficulty=media, tags={funzioni, dominio}]
Qual è il dominio della funzione $f(x)=\sqrt{\frac{x-1}{x+2}}$?

\begin{tabularx}{\linewidth}{@{}X X X X@{}}  
\boxed{A}\hspace{6pt} $(-\infty,-2)\cup[1,+\infty)$ &  
\boxed{B}\hspace{6pt} $[-2,1]$ &  
\boxed{C}\hspace{6pt} $(-2,1]$ &  
\boxed{D}\hspace{6pt} $(-\infty,-2]\cup[1,+\infty)$  
\end{tabularx}

\begin{solution}
\begin{enumerate}
    \item La risposta corretta è $A)$, perché deve essere $\frac{x-1}{x+2}\geq0$ e $x\neq-2$.
\end{enumerate}
\end{solution}
\end{question}

%----------------------------------------------------%

\begin{question}[ID=1.5, subtitle={Dominio di una funzione logaritmica}, difficulty=difficile, tags={funzioni, dominio}]
Qual è il dominio della funzione $f(x)=\ln\left(\frac{x-1}{x+1}\right)$?

\begin{tabularx}{\linewidth}{@{}X X X X@{}}  
\boxed{A}\hspace{6pt} $(-\infty,-1)\cup(1,+\infty)$ &  
\boxed{B}\hspace{6pt} $(-1,1)$ &  
\boxed{C}\hspace{6pt} $\mathbb{R}\setminus\{-1,1\}$ &  
\boxed{D}\hspace{6pt} $(-\infty,1)\cup(1,+\infty)$  
\end{tabularx}

\begin{solution}
\begin{enumerate}
    \item La risposta corretta è $A)$, perché l'argomento del logaritmo deve essere strettamente positivo: $\frac{x-1}{x+1}>0$.
\end{enumerate}
\end{solution}
\end{question}


%====================================================%
% SEZIONE 2: Segno di una funzione
%====================================================%

\begin{question}[ID=2.1, subtitle={Segno di una funzione lineare}, difficulty=facile, tags={funzioni, segno}]
Per quali valori di $x$ la funzione $f(x)=x-4$ è positiva?

\begin{tabularx}{\linewidth}{@{}X X X X@{}}  
\boxed{A}\hspace{6pt} $x<4$ &  
\boxed{B}\hspace{6pt} $x\leq4$ &  
\boxed{C}\hspace{6pt} $x>4$ &  
\boxed{D}\hspace{6pt} $x\geq4$  
\end{tabularx}

\begin{solution}
\begin{enumerate}
    \item La risposta corretta è $C)$, perché $x-4>0$ per $x>4$.
\end{enumerate}
\end{solution}
\end{question}

%----------------------------------------------------%

\begin{question}[ID=2.2, subtitle={Segno di una funzione quadratica}, difficulty=facile, tags={funzioni, segno}]
Per quali valori di $x$ la funzione $f(x)=x^2-9$ è negativa?

\begin{tabularx}{\linewidth}{@{}X X X X@{}}  
\boxed{A}\hspace{6pt} $x<-3\vee x>3$ &  
\boxed{B}\hspace{6pt} $-3<x<3$ &  
\boxed{C}\hspace{6pt} $x\leq-3\vee x\geq3$ &  
\boxed{D}\hspace{6pt} $-3\leq x\leq3$  
\end{tabularx}

\begin{solution}
\begin{enumerate}
    \item La risposta corretta è $B)$, perché $(x-3)(x+3)<0$ nell'intervallo compreso tra gli zeri.
\end{enumerate}
\end{solution}
\end{question}

%----------------------------------------------------%

\begin{question}[ID=2.3, subtitle={Segno di una funzione razionale}, difficulty=media, tags={funzioni, segno}]
Per quali valori di $x$ la funzione $f(x)=\frac{x+1}{x-2}$ è positiva?

\begin{tabularx}{\linewidth}{@{}X X X X@{}}  
\boxed{A}\hspace{6pt} $x<-1\vee x>2$ &  
\boxed{B}\hspace{6pt} $-1<x<2$ &  
\boxed{C}\hspace{6pt} $x<-2\vee x>1$ &  
\boxed{D}\hspace{6pt} $x>-1$  
\end{tabularx}

\begin{solution}
\begin{enumerate}
    \item La risposta corretta è $A)$, perché numeratore e denominatore hanno lo stesso segno per $x<-1$ oppure $x>2$.
\end{enumerate}
\end{solution}
\end{question}

%----------------------------------------------------%

\begin{question}[ID=2.4, subtitle={Segno di una funzione polinomiale}, difficulty=media, tags={funzioni, segno}]
Qual è l'insieme dei valori di $x$ per cui la funzione $f(x)=(x-1)(x+2)$ è positiva?

\begin{tabularx}{\linewidth}{@{}X X X X@{}}  
\boxed{A}\hspace{6pt} $-2<x<1$ &  
\boxed{B}\hspace{6pt} $x<-2\vee x>1$ &  
\boxed{C}\hspace{6pt} $x\leq-2\vee x\geq1$ &  
\boxed{D}\hspace{6pt} $x>1$  
\end{tabularx}

\begin{solution}
\begin{enumerate}
    \item La risposta corretta è $B)$, perché il prodotto è positivo all'esterno degli zeri $-2$ e $1$.
\end{enumerate}
\end{solution}
\end{question}

%----------------------------------------------------%

\begin{question}[ID=2.5, subtitle={Studio del segno di una funzione razionale}, difficulty=difficile, tags={funzioni, segno}]
Per quali valori di $x$ la funzione $f(x)=\frac{x^2-4}{x-1}$ è negativa?

\begin{tabularx}{\linewidth}{@{}X X X X@{}}  
\boxed{A}\hspace{6pt} $x<-2\vee1<x<2$ &  
\boxed{B}\hspace{6pt} $x<-2\vee x>2$ &  
\boxed{C}\hspace{6pt} $-2<x<1$ &  
\boxed{D}\hspace{6pt} $-2<x<1\vee x>2$  
\end{tabularx}

\begin{solution}
\begin{enumerate}
    \item La risposta corretta è $A)$, perché $\frac{(x-2)(x+2)}{x-1}<0$ negli intervalli $(-\infty,-2)$ e $(1,2)$.
\end{enumerate}
\end{solution}
\end{question}


%====================================================%
% SEZIONE 3: Parità di una funzione
%====================================================%

\begin{question}[ID=3.1, subtitle={Funzione pari}, difficulty=facile, tags={funzioni, parità}]
Quale delle seguenti funzioni è pari?

\begin{tabularx}{\linewidth}{@{}X X X X@{}}  
\boxed{A}\hspace{6pt} $f(x)=x^3+x$ &  
\boxed{B}\hspace{6pt} $f(x)=x^2+1$ &  
\boxed{C}\hspace{6pt} $f(x)=2x-1$ &  
\boxed{D}\hspace{6pt} $f(x)=x^3-2$  
\end{tabularx}

\begin{solution}
\begin{enumerate}
    \item La risposta corretta è $B)$, perché $f(-x)=(-x)^2+1=x^2+1=f(x)$.
\end{enumerate}
\end{solution}
\end{question}

%----------------------------------------------------%

\begin{question}[ID=3.2, subtitle={Funzione dispari}, difficulty=facile, tags={funzioni, parità}]
Quale delle seguenti funzioni è dispari?

\begin{tabularx}{\linewidth}{@{}X X X X@{}}  
\boxed{A}\hspace{6pt} $f(x)=x^2+1$ &  
\boxed{B}\hspace{6pt} $f(x)=x^3-2x$ &  
\boxed{C}\hspace{6pt} $f(x)=x^2-3$ &  
\boxed{D}\hspace{6pt} $f(x)=x+1$  
\end{tabularx}

\begin{solution}
\begin{enumerate}
    \item La risposta corretta è $B)$, perché $f(-x)=-x^3+2x=-f(x)$.
\end{enumerate}
\end{solution}
\end{question}

%----------------------------------------------------%

\begin{question}[ID=3.3, subtitle={Pari, dispari o nessuna delle due}, difficulty=media, tags={funzioni, parità}]
La funzione $f(x)=x^3+x^2$ è:

\begin{tabularx}{\linewidth}{@{}X X X X@{}}  
\boxed{A}\hspace{6pt} pari &  
\boxed{B}\hspace{6pt} dispari &  
\boxed{C}\hspace{6pt} né pari né dispari &  
\boxed{D}\hspace{6pt} costante  
\end{tabularx}

\begin{solution}
\begin{enumerate}
    \item La risposta corretta è $C)$, perché $f(-x)=-x^3+x^2$ non coincide né con $f(x)$ né con $-f(x)$.
\end{enumerate}
\end{solution}
\end{question}

%----------------------------------------------------%

\begin{question}[ID=3.4, subtitle={Parità di una funzione razionale}, difficulty=media, tags={funzioni, parità}]
La funzione $f(x)=\frac{1}{x}$ è:

\begin{tabularx}{\linewidth}{@{}X X X X@{}}  
\boxed{A}\hspace{6pt} pari &  
\boxed{B}\hspace{6pt} dispari &  
\boxed{C}\hspace{6pt} né pari né dispari &  
\boxed{D}\hspace{6pt} non è una funzione  
\end{tabularx}

\begin{solution}
\begin{enumerate}
    \item La risposta corretta è $B)$, perché $f(-x)=-\frac{1}{x}=-f(x)$.
\end{enumerate}
\end{solution}
\end{question}

%----------------------------------------------------%

\begin{question}[ID=3.5, subtitle={Parità di una funzione razionale}, difficulty=difficile, tags={funzioni, parità}]
Quale delle seguenti funzioni è pari?

\begin{tabularx}{\linewidth}{@{}X X X X@{}}  
\boxed{A}\hspace{6pt} $f(x)=\frac{x+1}{x^2+1}$ &  
\boxed{B}\hspace{6pt} $f(x)=\frac{x^2+1}{x^2-4}$ &  
\boxed{C}\hspace{6pt} $f(x)=\frac{x}{x^2-1}$ &  
\boxed{D}\hspace{6pt} $f(x)=\frac{x^3+1}{x^2+1}$  
\end{tabularx}

\begin{solution}
\begin{enumerate}
    \item La risposta corretta è $B)$, perché $f(-x)=\frac{x^2+1}{x^2-4}=f(x)$.
\end{enumerate}
\end{solution}
\end{question}


%====================================================%
% SEZIONE 4: Funzioni composte
%====================================================%

\begin{question}[ID=4.1, subtitle={Calcolo di una funzione composta}, difficulty=facile, tags={funzioni, composizione}]
Siano $f(x)=2x+1$ e $g(x)=x^2$. Quanto vale $(f\circ g)(3)$?

\begin{tabularx}{\linewidth}{@{}X X X X@{}}  
\boxed{A}\hspace{6pt} $7$ &  
\boxed{B}\hspace{6pt} $9$ &  
\boxed{C}\hspace{6pt} $19$ &  
\boxed{D}\hspace{6pt} $21$  
\end{tabularx}

\begin{solution}
\begin{enumerate}
    \item La risposta corretta è $C)$, perché $(f\circ g)(3)=f(g(3))=f(9)=19$.
\end{enumerate}
\end{solution}
\end{question}

%----------------------------------------------------%

\begin{question}[ID=4.2, subtitle={Espressione di una funzione composta}, difficulty=facile, tags={funzioni, composizione}]
Siano $f(x)=x+2$ e $g(x)=3x-1$. Qual è l'espressione di $(g\circ f)(x)$?

\begin{tabularx}{\linewidth}{@{}X X X X@{}}  
\boxed{A}\hspace{6pt} $3x+1$ &  
\boxed{B}\hspace{6pt} $3x+5$ &  
\boxed{C}\hspace{6pt} $x+5$ &  
\boxed{D}\hspace{6pt} $3x-3$  
\end{tabularx}

\begin{solution}
\begin{enumerate}
    \item La risposta corretta è $B)$, perché $(g\circ f)(x)=g(x+2)=3(x+2)-1=3x+5$.
\end{enumerate}
\end{solution}
\end{question}

%----------------------------------------------------%

\begin{question}[ID=4.3, subtitle={Composizione di funzioni algebriche}, difficulty=media, tags={funzioni, composizione}]
Siano $f(x)=x^2$ e $g(x)=x-1$. Qual è l'espressione di $(f\circ g)(x)$?

\begin{tabularx}{\linewidth}{@{}X X X X@{}}  
\boxed{A}\hspace{6pt} $x^2-1$ &  
\boxed{B}\hspace{6pt} $(x-1)^2$ &  
\boxed{C}\hspace{6pt} $x^2+1$ &  
\boxed{D}\hspace{6pt} $(x+1)^2$  
\end{tabularx}

\begin{solution}
\begin{enumerate}
    \item La risposta corretta è $B)$, perché $(f\circ g)(x)=f(x-1)=(x-1)^2$.
\end{enumerate}
\end{solution}
\end{question}

%----------------------------------------------------%

\begin{question}[ID=4.4, subtitle={Calcolo di una funzione composta con radicale}, difficulty=media, tags={funzioni, composizione}]
Siano $f(x)=\sqrt{x}$ e $g(x)=x+1$. Quanto vale $(f\circ g)(8)$?

\begin{tabularx}{\linewidth}{@{}X X X X@{}}  
\boxed{A}\hspace{6pt} $2$ &  
\boxed{B}\hspace{6pt} $3$ &  
\boxed{C}\hspace{6pt} $8$ &  
\boxed{D}\hspace{6pt} $9$  
\end{tabularx}

\begin{solution}
\begin{enumerate}
    \item La risposta corretta è $B)$, perché $(f\circ g)(8)=f(9)=\sqrt{9}=3$.
\end{enumerate}
\end{solution}
\end{question}

%----------------------------------------------------%

\begin{question}[ID=4.5, subtitle={Composizione di tre funzioni}, difficulty=difficile, tags={funzioni, composizione}]
Siano $f(x)=2x$, $g(x)=x+1$ e $h(x)=x^2$. Qual è l'espressione di $(f\circ g\circ h)(x)$?

\begin{tabularx}{\linewidth}{@{}X X X X@{}}  
\boxed{A}\hspace{6pt} $2x^2+1$ &  
\boxed{B}\hspace{6pt} $2x^2+2$ &  
\boxed{C}\hspace{6pt} $(2x+1)^2$ &  
\boxed{D}\hspace{6pt} $2(x+1)^2$  
\end{tabularx}

\begin{solution}
\begin{enumerate}
    \item La risposta corretta è $B)$, perché $(f\circ g\circ h)(x)=f(g(h(x)))=f(x^2+1)=2x^2+2$.
\end{enumerate}
\end{solution}
\end{question}